\newpage
\section*{TÓM TẮT}
\addcontentsline{toc}{section}{TÓM TẮT}

Công nghệ Thực tế tăng cường (Augmented Reality - AR) cho phép tích hợp các đối tượng ảo vào môi trường thực, tạo ra trải nghiệm tương tác trực quan giữa người dùng và hệ thống. Tuy nhiên, để các đối tượng ảo xuất hiện một cách tự nhiên và chân thực trong thế giới thực, hệ thống AR cần có khả năng hiểu và phân tích bối cảnh của môi trường xung quanh. Đây chính là mục tiêu của bài toán \textit{Scene Understanding \& Realism} trong AR.

Báo cáo này tập trung trình bày các khái niệm và kỹ thuật nền tảng giúp hệ thống AR hiểu cấu trúc không gian của môi trường và tăng mức độ chân thực khi hiển thị nội dung ảo. Cụ thể, báo cáo phân tích bốn thành phần kỹ thuật cốt lõi: ước lượng độ sâu và tái tạo hình học môi trường dưới dạng lưới 3D (depth estimation \& mesh reconstruction), kỹ thuật che khuất (occlusion) giữa vật thể thật và vật thể ảo, ước lượng ánh sáng (light estimation) và kết xuất bóng đổ (shadow rendering). Các kỹ thuật này cho phép các đối tượng ảo tương tác hợp lý với môi trường thực -- ví dụ như đứng trên bề mặt thật, bị che bởi vật thể thật, hoặc có ánh sáng và bóng đổ phù hợp với điều kiện chiếu sáng thực tế.

Thông qua việc tổng hợp và phân tích các kỹ thuật trên, báo cáo nhằm cung cấp cái nhìn tổng quan về vai trò của Scene Understanding trong việc nâng cao tính chân thực của hệ thống AR, đồng thời làm nền tảng cho việc thiết kế và triển khai các ứng dụng AR trong nhiều lĩnh vực khác nhau.

\newpage